% !TEX root = resume-en.tex

\documentclass[letterpaper,11pt]{article}

%----------LANGUAGE SWITCH----------
% This file holds both language versions; it is not compiled directly.
%   pdflatex resume-en  -> resume-en.pdf (English)
%   pdflatex resume-de  -> resume-de.pdf (Deutsch)
% The drivers select the language by defining \GERMANCV or not.
\newif\ifgerman
\ifdefined\GERMANCV\germantrue\else\germanfalse\fi

\usepackage{latexsym}
\usepackage[utf8]{inputenc}
\usepackage[T1]{fontenc}
\usepackage[empty]{fullpage}
\usepackage{titlesec}
\usepackage{marvosym}
\usepackage[usenames,dvipsnames]{color}
\usepackage{verbatim}
\usepackage{enumitem}
\usepackage[hidelinks]{hyperref}
\usepackage{fancyhdr}
\ifgerman
  \usepackage[ngerman]{babel}
\else
  \usepackage[english]{babel}
\fi
\usepackage{tabularx}
\input{glyphtounicode}


%----------FONT OPTIONS----------
% sans-serif
% \usepackage[sfdefault]{FiraSans}
% \usepackage[sfdefault]{roboto}
% \usepackage[sfdefault]{noto-sans}
% \usepackage[default]{sourcesanspro}

% serif
% \usepackage{CormorantGaramond}
% \usepackage{charter}


\pagestyle{fancy}
\fancyhf{} % clear all header and footer fields
\fancyfoot{}
\renewcommand{\headrulewidth}{0pt}
\renewcommand{\footrulewidth}{0pt}

% Adjust margins
\addtolength{\oddsidemargin}{-0.5in}
\addtolength{\evensidemargin}{-0.5in}
\addtolength{\textwidth}{1in}
\addtolength{\topmargin}{-.5in}
\addtolength{\textheight}{1.0in}

\urlstyle{same}

\raggedbottom
\raggedright
\setlength{\tabcolsep}{0in}

% Sections formatting
\titleformat{\section}{
  \vspace{-4pt}\scshape\raggedright\large
}{}{0em}{}[\color{black}\titlerule \vspace{-5pt}]

% Ensure that generate pdf is machine readable/ATS parsable
\pdfgentounicode=1

%-------------------------
% Custom commands

% \bi{English}{Deutsch} -- picks the language selected by the switch above
\newcommand{\bi}[2]{\ifgerman#2\else#1\fi}

\newcommand{\resumeItem}[1]{
  \item\small{
    {#1 \vspace{-2pt}}
  }
}

\newcommand{\resumeSubheading}[4]{
  \vspace{-2pt}\item
    \begin{tabular*}{0.97\textwidth}[t]{l@{\extracolsep{\fill}}r}
      \textbf{#1} & #2 \\
      \textit{\small#3} & \textit{\small #4} \\
    \end{tabular*}\vspace{-7pt}
}

\newcommand{\resumeSubSubheading}[2]{
    \item
    \begin{tabular*}{0.97\textwidth}{l@{\extracolsep{\fill}}r}
      \textit{\small#1} & \textit{\small #2} \\
    \end{tabular*}\vspace{-7pt}
}

\newcommand{\resumeProjectHeading}[2]{
    \item
    \begin{tabular*}{0.97\textwidth}{l@{\extracolsep{\fill}}r}
      \small#1 & #2 \\
    \end{tabular*}\vspace{-7pt}
}

\newcommand{\resumeSubItem}[1]{\resumeItem{#1}\vspace{-4pt}}

\renewcommand\labelitemii{$\vcenter{\hbox{\tiny$\bullet$}}$}

\newcommand{\resumeSubHeadingListStart}{\begin{itemize}[leftmargin=0.15in, label={}]}
\newcommand{\resumeSubHeadingListEnd}{\end{itemize}}
\newcommand{\resumeItemListStart}{\begin{itemize}}
\newcommand{\resumeItemListEnd}{\end{itemize}\vspace{-5pt}}

%-------------------------------------------
%%%%%%  RESUME STARTS HERE  %%%%%%%%%%%%%%%%%%%%%%%%%%%%


\begin{document}

%----------HEADING----------
% \begin{tabular*}{\textwidth}{l@{\extracolsep{\fill}}r}
%   \textbf{\href{http://sourabhbajaj.com/}{\Large Sourabh Bajaj}} & Email : \href{mailto:sourabh@sourabhbajaj.com}{sourabh@sourabhbajaj.com}\\
%   \href{http://sourabhbajaj.com/}{http://www.sourabhbajaj.com} & Mobile : +1-123-456-7890 \\
% \end{tabular*}

\begin{center}
    \textbf{\Huge \scshape Jonas Christian Öchsner} \\ \vspace{3pt}
    \small Schleifweg 23, 97273 Kürnach, \bi{Germany}{Deutschland} \\ \vspace{2pt}
    \small +49 1573 5544890 $|$ \href{mailto:jonasoechsner06@gmail.com}{\underline{jonasoechsner06@gmail.com}} $|$
    \href{mailto:j_oechsner@stud.uni-bamberg.de}{\underline{j\_oechsner@stud.uni-bamberg.de}}
\end{center}


%-----------EDUCATION-----------
\section{\bi{Education}{Ausbildung}}
  \resumeSubHeadingListStart
    \resumeSubheading
      {TU Dublin}{Dublin, \bi{Ireland}{Irland}}
      {\bi{ERASMUS Exchange (Computing)}{ERASMUS-Auslandsstudium (Informatik)}}{\bi{Sept. 2026 -- May 2027}{Sep. 2026 -- Mai 2027}}
    \resumeSubheading
      {Universität Bamberg}{Bamberg, \bi{Germany}{Deutschland}}
      {\bi{B.Sc. Angewandte Informatik (Applied Computer Science)}{B.Sc. Angewandte Informatik}}{\bi{Oct. 2024 -- PRESENT}{Okt. 2024 -- HEUTE}}
    \resumeSubheading
      {Röntgen-Gymnasium Würzburg}{Würzburg, \bi{Germany}{Deutschland}}
      {\bi{Abitur (Grade 2.0)}{Abitur (Note 2,0)}}{\bi{Sep. 2016 -- Jul. 2024}{Sep. 2016 -- Juli 2024}}
  \resumeSubHeadingListEnd


%-----------EXPERIENCE-----------
\section{\bi{Experience}{Berufserfahrung}}
  \resumeSubHeadingListStart

    \resumeSubheading
      {\bi{Software Development for AI-Powered Data Protection Classification}{Softwareentwicklung für KI-gestützte Datenschutzklassifikation}}{\bi{Oct. 2025 -- PRESENT}{Okt. 2025 -- HEUTE}}
      {LIfBi: Leibniz-Institut für Bildungsverläufe}{Bamberg, \bi{Germany}{Deutschland}}
      \resumeItemListStart
        \resumeItem{\bi
          {Run an own agile-managed software project (see below), from requirements engineering through maintenance}
          {Eigenverantwortliche Leitung eines agil gesteuerten Softwareprojekts (siehe unten), von der Anforderungserhebung bis zur Wartung}}
        \resumeItem{\bi
          {Ensure the solution complies with European data protection law}
          {Sicherstellung der Konformität der Lösung mit dem europäischen Datenschutzrecht}}
        \resumeItem{\bi
          {Handle customer communication and set up Continuous Integration for the project}
          {Kundenkommunikation sowie Aufbau der Continuous Integration für das Projekt}}
      \resumeItemListEnd

    \resumeSubheading
      {\bi{Tutor, ``Introduction to AI''}{Tutor, „Einführung in die KI“}}{\bi{Jul. 2026 -- PRESENT}{Juli 2026 -- HEUTE}}
      {\bi{Cognitive Systems Chair, Otto-Friedrich University Bamberg}{Lehrstuhl für Kognitive Systeme, Otto-Friedrich-Universität Bamberg}}{Bamberg, \bi{Germany}{Deutschland}}
      \resumeItemListStart
        \resumeItem{\bi
          {Teach and support students in the ``Introduction to AI'' course}
          {Lehre und Betreuung der Studierenden im Kurs „Einführung in die KI“}}
      \resumeItemListEnd

    \resumeSubheading
      {\bi{Temporary Delivery Driver}{Aushilfsfahrer (Auslieferung)}}{\bi{Jul. 2024 -- Aug. 2024}{Juli 2024 -- Aug. 2024}}
      {Viktor Nußbaumer Bestes für Küche und Gastlichkeit GmbH \& Co. KG}{Kürnach, \bi{Germany}{Deutschland}}

    \resumeSubheading
      {\bi{Picker}{Kommissionierer}}{\bi{Jul. 2022 -- Aug. 2022}{Juli 2022 -- Aug. 2022}}
      {Hans Geis GmbH + Co KG}{Kürnach, \bi{Germany}{Deutschland}}

  \resumeSubHeadingListEnd


%-----------PROJECTS-----------
\section{\bi{Projects (both closed source)}{Projekte (beide Closed Source)}}
    \resumeSubHeadingListStart
      \resumeProjectHeading
          {\textbf{\bi{Raven's Progressive Matrices -- Abstract Visual Reasoning}{Raven's Progressive Matrices -- abstraktes visuelles Schließen}} $|$ \emph{Python, PyTorch}}{\bi{Mar. 2026 -- Aug. 2026}{März 2026 -- Aug. 2026}}
          \resumeItemListStart
            \resumeItem{\bi
              {Team project (Cognitive Systems, Universität Bamberg) on trying to replicate ``An Interpretable Neuro-symbolic Model for Raven's Progressive Matrices Reasoning'', Zhao et al. (2023)}
              {Teamprojekt (Kognitive Systeme, Universität Bamberg) zur Replikation von „An Interpretable Neuro-symbolic Model for Raven's Progressive Matrices Reasoning“, Zhao et al. (2023)}}
            \resumeItem{\bi
              {Reproduced two published dataset generators (RAVEN, RAVEN-FAIR)}
              {Nachbildung zweier veröffentlichter Datensatzgeneratoren (RAVEN, RAVEN-FAIR)}}
            \resumeItem{\bi
              {Reproduced the three-stage neuro-symbolic reasoning pipeline (CNN $\rightarrow$ supervised VAE for semantic attribute extraction $\rightarrow$ cognitive-map relational reasoning) from the paper, achieving 94\% Prediction Accuracy}
              {Reproduktion der dreistufigen neuro-symbolischen Reasoning-Pipeline aus dem Paper (CNN $\rightarrow$ überwachter VAE zur semantischen Attributextraktion $\rightarrow$ relationales Schlussfolgern auf einer Cognitive Map) mit 94\% Vorhersagegenauigkeit}}
          \resumeItemListEnd
      \resumeProjectHeading
          {\textbf{\bi{LLM-Powered GDPR Compliance Classifier}{LLM-basierter DSGVO-Compliance-Klassifikator}} $|$ \emph{Python, PostgreSQL, Docker, Ollama}}{\bi{Oct. 2025 -- PRESENT}{Okt. 2025 -- HEUTE}}
          \resumeItemListStart
            \resumeItem{\bi
              {Built a tool for LIfBi that classifies survey questions against the special categories of personal data from Art. 9 GDPR}
              {Entwicklung eines Tools für das LIfBi, das Befragungsfragen anhand der besonderen Kategorien personenbezogener Daten nach Art. 9 DSGVO klassifiziert}}
            \resumeItem{\bi
              {Developed full stack: SvelteKit/TypeScript frontend and API routes, PostgreSQL schema, and a three-container Docker Compose deployment (load balancing with Kubernetes is planned)}
              {Umsetzung des gesamten Stacks: SvelteKit/TypeScript-Frontend und API-Routen, PostgreSQL-Schema sowie ein Deployment aus drei Containern mit Docker Compose (Load balancing mit Kubernetes ist geplant)}}
            \resumeItem{\bi
              {Ran all LLM inference locally through Ollama to protect sensitive institute data}
              {Lokale LLM-Inferenz über Ollama zum Schutz sensibler Institutsdaten}}
          \resumeItemListEnd
    \resumeSubHeadingListEnd


%-----------HIGHLIGHTED GRADES-----------
\section{\bi{Highlighted Grades (Average Score 1.98}{Ausgewählte Noten (Notendurchschnitt 1,98}, \bi{99 credited ETCS}{99 angerechnete ETCS})}
 \begin{itemize}[leftmargin=0.15in, label={}]
    \small{\item{
     \bi
      {Statistics \& Data Science (2.0) $|$ Discrete Modelling (1.3) $|$ Java Programming (1.3) $|$ Linear Algebra (1.3) \\
       DevOps for AI Systems (1.3) $|$ Analysis (1.7) $|$ Introduction to AI (1.7) $|$ Introduction to ML (1.7) $|$ European Data Protection Law (2.3)}
      {Statistik \& Data Science (2,0) $|$ Diskrete Modellierung (1,3) $|$ Java-Programmierung (1,3) $|$ Lineare Algebra (1,3) \\
       DevOps für KI-Systeme (1,3) $|$ Analysis (1,7) $|$ Einführung in die KI (1,7) $|$ Einführung in ML (1,7), $|$ Europäisches Datenschutzrecht (2,3)}
    }}
 \end{itemize}


%-----------SKILLS-----------
\section{\bi{Skills}{Kenntnisse}}
 \begin{itemize}[leftmargin=0.15in, label={}]
    \small{\item{
     \textbf{\bi{Languages}{Programmiersprachen}}{: \bi
       {Python, Java, R, SPSS (introductory), C-style languages (introductory)}
       {Python, Java, R, SPSS (Grundkenntnisse), C-artige Sprachen (Grundkenntnisse)}} \\
     \textbf{\bi{Libraries \& Tools}{Bibliotheken \& Werkzeuge}}{: PyTorch, NumPy, Pandas, UML, Git} \\
     \textbf{\bi{Concepts}{Konzepte}}{: \bi
       {Machine Learning, Data Science, Distributed Systems, Object-Oriented Architecture Patterns}
       {Maschinelles Lernen, Data Science, Verteilte Systeme, objektorientierte Architekturmuster}} \\
     \textbf{\bi{Practices}{Methoden}}{: \bi
       {Agile Project Management, Requirements Engineering, Continuous Integration, Software Maintenance}
       {Agiles Projektmanagement, Requirements Engineering, Continuous Integration, Softwarewartung}} \\
     \textbf{\bi{Spoken Languages}{Sprachen}}{: \bi
       {German (native), English (certified B2)}
       {Deutsch (Muttersprache), Englisch (zertifiziert B2)}} \\
    }}
 \end{itemize}


%-------------------------------------------
\end{document}
